\section{合成路线与相控制}

\subsection{固相与溶液前驱体}

常规固相反应路线以 BaCO$_3$\slash BaO、ZnO（或 ZnCO$_3$）和 SiO$_2$ 为原料，经混合、预烧、研磨和高温保温制得 BaZn$_2$Si$_2$O$_7$；Co、Mg 等取代体系通常沿用同一氧化物/碳酸盐反应路径，并通过中间研磨缩短扩散距离 \cite{kerstan2013bazn2si2o7,thieme2016new}。燃烧辅助法可减小前驱体的混合尺度，但 Eu$^{2+}$ 发光材料的研究表明，最终相组成仍取决于后续煅烧温度和气氛 \cite{yao2010synthesis,yao2010synthesisx}。溶胶--凝胶路线在分子尺度预混合 Ba、Sr、Zn 和 Si，适合制备细粉和致密陶瓷；Kracker 等借此实现了高温和低温多晶型的可控制备 \cite{kracker2016sol}。

\subsection{玻璃析晶路线}

BaO\slash SrO\slash ZnO\slash SiO$_2$ 玻璃可经成核--晶化热处理析出 Ba$_{1-x}$Sr$_x$Zn$_2$Si$_2$O$_7$。原位显微观察和 HT\nobreakdash-XRD 表明，Pt、Sb$_2$O$_3$、Nb$_2$O$_5$、Ta$_2$O$_5$、P$_2$O$_5$ 等成核剂会改变表面析晶与体析晶的竞争关系 \cite{bocker2012in,vladislavova2017heterogeneous,vladislavova2019crystallization,heitmann2019surface,thieme2017on,vladislavova2018nucleation}。ZrO$_2$\slash TiO$_2$ 和 Al$_2$O$_3$ 添加剂则同时影响黏度、成核密度及最终晶粒尺寸 \cite{vladislavova2018crystallisation,thieme2017effect}。对于 BaZn$_2$Si$_2$O$_7$ 本身，玻璃路线的优势在于成分均匀和晶粒尺寸可调，不足之处是残余玻璃相与微裂纹会掩盖本征热膨胀 \cite{thieme2016thermal,thieme2015high}。表~\ref{tab:routes}\CJKecglue{}对照了各路线的前驱体、主相判据与局限。

\begin{table}[bp]
\centering
\caption{已报道合成路线的条件--结果对照。}
\label{tab:routes}
\begin{tabular}{P{\dimexpr 0.1571\textwidth-2\tabcolsep\relax}P{\dimexpr 0.3049\textwidth-2\tabcolsep\relax}P{\dimexpr 0.2373\textwidth-2\tabcolsep\relax}P{\dimexpr 0.3007\textwidth-2\tabcolsep\relax}}
\toprule
    \thead{路线} & \thead{前驱体/关键操作} & \thead{主相判据} & \thead{局限}\\
\midrule
    混合氧化物固相反应 & 氧化物或碳酸盐原料；多次研磨与高温保温 & 粉末 XRD 与热膨胀结果一致 & 扩散受限；易残留 BaZnSiO$_4$ 等伴生相 \cite{kerstan2013bazn2si2o7}\\
    \addlinespace[3pt]
    溶胶--凝胶/燃烧合成 & 分子级混合；凝胶干燥或燃烧后煅烧 & 粒径小；可获得多晶型粉体 & 有机残余物和气氛影响价态 \cite{kracker2016sol,yao2010synthesis}\\
    \addlinespace[3pt]
    玻璃析晶 & BaO\slash SrO\slash ZnO\slash SiO$_2$ 玻璃；成核与晶化热处理 & 原位 XRD/显微观察确定晶化前沿 & 残余玻璃相与微裂纹；成核剂引入杂相 \cite{bocker2012in,vladislavova2017heterogeneous}\\
\bottomrule
\end{tabular}
\end{table}

\subsection{相图与热力学数据}

现有文献已对 BaO--ZnO 二元体系进行相平衡评估，但未检索到可直接用于 BaO--ZnO--SiO$_2$ 三元体系的完整 CALPHAD（相图计算）数据库。因此，组成点是否落在 BaZn$_2$Si$_2$O$_7$ 初晶区、是否经历共晶液相，只能结合 XRD 实验和热分析判断。负热膨胀硅酸盐的理论--实验研究表明，声子自由能和温度相关的相稳定性不能仅由室温结构推断 \cite{erlebach2021phase}。
